\chapter*{Learning 与硬件资料融合索引}
\addcontentsline{toc}{chapter}{Learning 与硬件资料融合索引}
\markboth{Learning 与硬件资料融合索引}{Learning 与硬件资料融合索引}

本索引列出每一组 learning 材料在书中的落点。书稿按因果依赖重组内容；
读者可以顺序读书，也可以从原学习分册或硬件图反查对应章节。

\newcommand{\learningpathii}[2]{%
  \shortstack[l]{\scriptsize\filepath{#1}\\[-0.15em]\scriptsize\filepath{#2}}}
\newcommand{\learningpathiii}[3]{%
  \shortstack[l]{\scriptsize\filepath{#1}\\[-0.15em]\scriptsize\filepath{#2}\\[-0.15em]
  \scriptsize\filepath{#3}}}

\section*{一、融合时采用的五条规则}

\begin{enumerate}[leftmargin=2.2em]
  \item \textbf{当前实现优先。}若早期 v0.x 分册与当前 v2.0 在容量、数据结构
        或路径上不同，正文先解释历史基线，再明确当前增量，不把旧数字冒充现状。
  \item \textbf{背景知识前置。}第一次使用 ROM、RAM、地址、寄存器、端口、
        中断、页表、进程、fd 等概念时先解释“是什么和为什么”，再进入项目细节。
  \item \textbf{同一因果链合并。}例如键盘硬件、IRQ、FIFO、阻塞、fd 0 和 Shell
        来自多个分册，在书中合并成一条端到端路径，避免读者自行拼接。
  \item \textbf{QEMU 与实体板分栏。}QEMU 图证明架构/设备事务，实体原理图证明
        电源和铜线；共同概念相互对照，平台专属假设明确标注。
  \item \textbf{图与文字互相校验。}图中每个主要节点在正文都有输入、输出、
        初始化者、完成条件和失败现象；正文声称的完整路径也必须能在图上追踪。
\end{enumerate}

\section*{二、零基础四章怎样吸收全仓资料}

\begin{osgridlongtable}{L{0.12\linewidth}L{0.37\linewidth}L{0.41\linewidth}}
章节 & 主要本地事实来源 & 重组后的教学职责 \\ \hline
\endfirsthead
章节 & 主要本地事实来源 & 重组后的教学职责 \\ \hline
\endhead
第 1 章 &
\learningpathiii{docs/learning/background/}{01-computer-pc-and-x86-}{history.md}、
\learningpathiii{docs/learning/background/}{02-data-address-assembly-}{and-abi.md}，
以及项目词汇表与三档容量规格 &
从量、单位、进制、状态和编码建立计算机语言，并把规格数字变成可检查的
范围、单位和数量级。\\ \hline
第 2 章 &
\learningpathiii{docs/learning/}{physical-carrier-}{circuit-guide.md}、实体原理图、
manufacturing checklist 与初学者硬件词汇表 &
把电源、地、电阻、电容、二极管、MOSFET、上拉、去耦、RC、符号、网络、
封装和实物方向连成可计算电路。\\ \hline
第 3 章 &
\learningpathiii{docs/learning/background/}{01-computer-pc-and-x86-}{history.md}、
\learningpathii{hardware-assembly-}{and-wiring.md} 与硬件状态图 &
用 CMOS、逻辑门、全加器、MUX、译码器、触发器、setup/hold、亚稳态和 FSM
解释处理器与设备控制器的共同数字积木。\\ \hline
第 4 章 &
\filepath{docs/background_knowledge.md}、\filepath{docs/architecture.md}、
\filepath{docs/modules/firmware.md}、\filepath{docs/modules/stage1.md}、
整机/启动图与实体载板资料 &
把芯片封装和引脚、ISA/微架构、CPU 数据通路、总线时序、SRAM/DRAM/Flash、
地址译码和复位取指接成一台完整机器。\\ \hline
\end{osgridlongtable}

\section*{三、总览、前置与背景分册怎样进入正文}

\begin{osgridlongtable}{L{0.25\linewidth}L{0.33\linewidth}L{0.32\linewidth}}
Learning 来源 & 核心知识链 & 成书中的主要落点 \\ \hline
\endfirsthead
Learning 来源 & 核心知识链 & 成书中的主要落点 \\ \hline
\endhead
\filepath{docs/learning/README.md} &
项目边界、学习顺序、四遍阅读法、构建运行、生产源码覆盖与完成判定 &
前言与阅读地图；第 5 章机器边界；第 31 章实践与检查 \\ \hline
\learningpathii{00-prerequisites-}{and-hardware.md} &
复位、四类地址、数据表示、x86 模式、寄存器、描述符、中断、分页、设备、
磁盘/内存、ELF、freestanding、QEMU/GDB &
第 5--12 章的基础主题；零基础机器章；初学者术语表 \\ \hline
\filepath{background/README.md} &
背景专题地图、统一符号、反事实阅读方法和通过标准 &
前言中的阅读方法；第 5 章软件边界；书末检查清单 \\ \hline
\learningpathiii{background/01-computer-}{pc-and-x86-}{history.md} &
计算机角色、取指循环、ISA/微架构/平台、电源复位、8086--AMD64、PC 设备史 &
第 1--5 章机器部件与软件边界；第 7 章 x86 四十多年兼容史；
第 8、11 章平台与模式切换 \\ \hline
\learningpathiii{background/02-data-}{address-assembly-}{and-abi.md} &
bit、整数、溢出、字节序、对齐、地址、汇编、栈、函数/系统调用 ABI &
第 1--4 章零基础；第 6 章数据与汇编；第 9、11、17--19 章 ABI 边界 \\ \hline
\learningpathiii{background/03-toolchain-}{elf-and-freestanding-}{cpp.md} &
host/target、预处理到链接、ELF、ROM/raw image、freestanding C++、产物审计 &
第 9 章工具链；第 11 章 ELF loader；第 31 章修改与验证方法 \\ \hline
\learningpathiii{background/04-memory-}{paging-protection-}{and-exceptions.md} &
RAM/cache、E820、frame、四级分页、权限、TLB、GDT/TSS/IDT、异常、用户指针 &
第 12--16 章异常与内存；第 17 章用户边界；第 23 章虚拟内存；术语表 \\ \hline
\learningpathiii{background/05-interrupts-}{time-and-pc-}{devices.md} &
Port I/O/MMIO、polling/timeout、PIC/LAPIC、PIT、UART、i8042、ATA、QMP &
第 8 章 Port I/O 全图；第 10 章 UART 与测试；第 17 章中断/键盘全图 \\ \hline
\learningpathiii{background/06-processes-}{syscalls-files-}{and-shell.md} &
Kernel/User、进程/线程、调度、系统调用、阻塞、锁、pipe、fd、文件系统、Shell &
第 17 章特权边界；第 18--30 章进程、对象、VM、持久化与用户环境 \\ \hline
\learningpathiii{background/07-testing-}{debugging-and-}{evidence.md} &
oracle、分层测试、随机模型、产物审计、QEMU、故障注入、持久性、调试 &
第 10 章观测与测试；各章失败框；第 31 章实践；书末实验与检查清单 \\ \hline
\end{osgridlongtable}

\section*{四、历史阶段分册怎样进入三十一章}

\begin{osgridlongtable}{L{0.25\linewidth}L{0.34\linewidth}L{0.31\linewidth}}
阶段分册 & 被保留的完整机制 & 成书落点 \\ \hline
\endfirsthead
阶段分册 & 被保留的完整机制 & 成书落点 \\ \hline
\endhead
\learningpathii{01-v0.0-engineering-}{baseline.md} &
CMake/Preset、host/target、AddressRange、测试分层、失败语义、源码顺序 &
第 5 章软件边界；第 9 章工具链；第 10/31 章测试与实践 \\ \hline
\learningpathii{02-v0.1-reset-}{and-serial.md} &
复位状态、段/栈、COM1、PIT、ROM 布局、串口协议、故障注入 &
第 7 章复位；第 9 章 ROM；第 10 章 UART；第 11 章启动全图 \\ \hline
\learningpathii{03-v0.2-firmware-}{stage1-loader.md} &
Stage 1 磁盘描述符、RAM 窗口、ATA PIO、校验、远转移、镜像生成 &
第 11 章 loader 主线与启动/内存全图 \\ \hline
\learningpathii{04-v0.3-}{long-mode.md} &
A20、GDT、保护模式、临时页表、PAE/LME/PG、64 位入口与故障标记 &
第 7 章历史依赖；第 11 章模式转换及失败定位 \\ \hline
\learningpathii{05-v0.4-kernel-elf-}{and-boot-info.md} &
Kernel 容器、CRC32、ELF64、严格 PT\_LOAD、两遍装载、BSS、BootInfo、C++ &
第 9 章 ELF/freestanding；第 11 章装载与交接 \\ \hline
\learningpathiii{06-v0.5-descriptor-}{tables-and-}{exceptions.md} &
GDT/TSS/IST/IDT、异常 frame、汇编 adapter、panic、自检和三重故障 &
第 12 章描述符、异常与内核接管 \\ \hline
\learningpathii{07-v0.6-memory-}{management.md} &
E820、物理 frame、正式页表、NX/WP、guard、堆及当前可回收内存演进 &
第 13--16 章物理页、页表、KVA、堆与资源管理 \\ \hline
\learningpathii{08-v0.7-interrupts-}{and-devices.md} &
PIC/LAPIC、PIT、PS/2、Kernel ATA、初始化顺序、中断安全和设备测试 &
第 17 章中断/设备主线及两张整页矢量图 \\ \hline
\learningpathii{09-v0.8-ring3-}{and-system-calls.md} &
用户 ELF、四道 Ring 3 门、IRETQ、INT 0x80、用户复制与异常隔离 &
第 17 章用户边界；第 19 章系统调用入口 \\ \hline
\learningpathii{10-v0.9-processes-}{and-scheduling.md} &
PCB、独立 CR3/Ring 0 栈、176 B frame、PIT 抢占、退出和资源守恒 &
第 18 章进程/调度；第 15 章动态栈与页表回收 \\ \hline
\learningpathii{11-v0.10-synchronization-}{and-ipc.md} &
SpinLock、Blocked/Ready、定向唤醒、64 B 管道、丢失唤醒与用户复制 &
第 18、28 章同步/IPC；键盘到 Shell 的阻塞链 \\ \hline
\learningpathii{12-v0.11-}{file-system.md} &
磁盘所有权、superblock、inode、bitmap、目录、LRU cache、事务和跨启动 &
第 24 章文件系统盘面；第 26--27 章写回与恢复 \\ \hline
\learningpathii{13-v1.0-user-}{environment.md} &
i8042、Console FIFO、统一 fd、Try/Wait、idle、Ring 3 Shell 和自动输入 &
第 17 章按键全链；第 30 章用户环境；第 31 章整体验收 \\ \hline
\learningpathii{14-v1.6-rootfs-}{v2.md} &
rootfs v2 盘面、三级间接树、稀疏文件、完整命名空间修改、fsck 与跨启动 &
第 25 章 rootfs v2 与大文件边界 \\ \hline
\learningpathii{15-v1.7-pid1-process-}{tree-exec.md} &
PID1、父子树、Zombie/孤儿、磁盘 ELF、参数栈、spawn/exec/wait 与事务回滚 &
第 22 章 PID 1、进程树与程序映像 \\ \hline
\learningpathii{16-v1.8-anonymous-vma-}{demand-paging.md} &
VMA/PTE 分层、匿名 reservation、按需缺页、连续栈、brk、用户堆与资源守恒 &
第 23 章匿名 VMA、用户堆与虚拟内存演进 \\ \hline
\learningpathiii{17-v1.9-file-backed-}{vma-lazy-elf-}{page-cache.md} &
稳定文件身份、FileBacking、按需 ELF、clean cache、shared/private 与失效 &
第 23 章文件后备 VMA、按需 ELF 与页缓存 \\ \hline
\learningpathii{18-v1.10-fork-copy-}{on-write.md} &
fork 返回、private COW、CopyToUser、fd/cwd/后备继承和两阶段失败回滚 &
第 23 章 fork/COW 与多线程 fork/exec 边界 \\ \hline
\learningpathiii{19-v1.11-unix-io-}{external-}{shell.md} &
动态 Pipe、pipe/dup2、外部工具、重定向、16 级流水线与跨层资源守恒 &
第 28 章 Unix I/O；第 30 章外部工具与用户环境 \\ \hline
\learningpathiii{20-v1.12-user-threads-}{tls-private-}{futex.md} &
用户 Thread、FS-base TLS、private futex、Mutex/CV/Once、join 与多线程 exec &
第 28 章用户线程、futex 与并发演进 \\ \hline
\learningpathiii{21-v1.13-monotonic-clock-}{deadline-timed-}{wait.md} &
PIT 有理数换算、单调纳秒、deadline queue、sleep、timed futex 与单赢家竞争 &
第 28 章单调时间、deadline 与可中断等待 \\ \hline
\learningpathiii{22-v1.14-process-}{signals-}{sigreturn.md} &
处置、mask、pending、进程组、单赢家唤醒、用户 frame 与严格 sigreturn &
第 29 章进程信号、sigreturn 与作业控制前置 \\ \hline
\learningpathiii{23-v1.15-tty-}{session-job-}{control.md} &
行规、控制字符、控制终端、session/进程组、停止/继续事件与交互式作业控制 &
第 29 章 TTY、session 与交互式作业控制 \\ \hline
\learningpathiii{24-v1.16-irq14-}{block-request-}{writeback.md} &
ATA/PIC 标志、单飞请求、IRQ/timeout 单赢家、BlockIo、写通知与页状态写回 &
第 26 章 IRQ14、BlockIo 与 writeback \\ \hline
\learningpathiii{25-v1.17-ordered-}{metadata-}{journal.md} &
WAL 历史、credit、盘面记录、FLUSH 顺序、断电矩阵、checkpoint 与幂等 replay &
第 27 章 ordered metadata journal 与恢复边界 \\ \hline
\learningpathiii{26-v1.18-abi-v2-}{devfs-procfs-}{release-freeze.md} &
ABI 版本/offset、ELF 身份、devfs 分层、procfs 快照与 32 个用户工具 &
第 30 章 ABI/devfs/procfs \\ \hline
\learningpathiii{27-v2.0-}{integration-}{release.md} &
三档回归、跨层代码走读、故障定位与后续演进边界 &
第 31 章系统回顾、测试方法与下一步路线 \\ \hline
\end{osgridlongtable}

\section*{五、硬件图册与实体资料逐项落点}

\begin{osgridlongtable}{L{0.29\linewidth}L{0.27\linewidth}L{0.34\linewidth}}
来源 & 图/资料回答的问题 & 成书呈现 \\ \hline
\endfirsthead
来源 & 图/资料回答的问题 & 成书呈现 \\ \hline
\endhead
\learningpathii{hardware-assembly-}{and-wiring.md} &
十三阶段怎样装配；整机、启动/内存、Port I/O、中断、按键和持久化路径 &
第 4、8、11、17、24--30 章的横置整页矢量图及逐图讲解 \\ \hline
\filepath{assets/hardware_stage_roadmap.svg} &
v0.0--v1.0 十三个能力增量及 v1.1--v1.8 延伸 &
图~\ref{fig:hardware-stage-roadmap} 与四阶段依赖表 \\ \hline
\filepath{assets/x86\_64\_os\_hardware\_wiring.svg} &
Host、ROM、RAM、CPU、Port I/O、PIC/LAPIC 与设备的全局关系 &
图~\ref{fig:system-wiring-overview} 与四种线/五块硬件解释 \\ \hline
\filepath{assets/boot\_and\_memory\_wiring.svg} &
五次控制权交接、三种 CPU 模式和关键物理区间 &
图~\ref{fig:boot-memory-wiring} 与所有权/故障定位 \\ \hline
\filepath{assets/port\_io\_topology.svg} &
六组端口、访问宽度、状态协议及 MMIO 边界 &
图~\ref{fig:port-io-topology} 与逐设备寄存器解释 \\ \hline
\filepath{assets/interrupt\_routing.svg} &
IRQ0/1、PIC、LAPIC、IDT、handler、EOI 和 IRETQ &
图~\ref{fig:interrupt-routing-full} 与硬件/软件泳道 \\ \hline
\filepath{assets/keyboard\_to\_shell.svg} &
QMP 到 PS/2、IRQ、FIFO、fd 0、Shell 的十四步链 &
图~\ref{fig:keyboard-to-shell-full} 与逐步失败矩阵 \\ \hline
\filepath{assets/storage\_persistence.svg} &
Ring 3 写入、文件对象、cache、ATA、raw disk 与新启动验证 &
图~\ref{fig:storage-persistence-full} 与持久性证据层级 \\ \hline
\learningpathii{physical-carrier-}{circuit-guide.md} &
Mu/DFR1142 电源、高速、低速、260-pin、PCB 与实机适配 &
第 8 章“实体 x86-64 载板全链路案例” \\ \hline
三张 \learningpathii{assets/physical\_}{carrier/*.svg} &
电源、高速和控制接口逐引脚/逐网络简化图 &
第 8 章三张矢量学习图；377 个显式引脚与 210 段导线门禁 \\ \hline
DFR1142 十页上游原理图 PDF/KiCad &
全部器件位号、数值、网络、NC/DNP 与 PCB 事实 &
第 8 章按功能嵌入十页原始矢量原理图并逐页解释 \\ \hline
Mu pinout 与 260-pin 表 &
模块正反方向、辅助连接器、奇偶触点和每个信号定义 &
纯 TikZ 矢量方向图、关键 pin 表与 HSIO 分配表 \\ \hline
\learningpathii{hardware/.../}{provenance.md} &
上游提交、许可、哈希与事实边界 &
实体案例来源段、图片 README 与规范/延伸阅读 \\ \hline
\learningpathii{hardware/.../}{manufacturing\_checklist.md} &
ERC/DRC、BOM、叠层、阻抗、热、机械和上电前检查 &
实体案例 bring-up 段与书末硬件验收清单 \\ \hline
\end{osgridlongtable}

\section*{六、为什么不把同一段文字复制两遍}

Learning 分册保留按版本实验的阅读顺序，书稿则保留按机制依赖的长篇叙事。
例如 \code{TSS.RSP0} 同时出现在 v0.5 异常、v0.8 Ring 3、v0.9 进程和 v1.3
SYSCALL 分册中；书中先在处理器契约处解释 TSS，再分别说明异常换栈、每进程
Ring 0 栈和 SYSCALL 可信换栈，不把四段近似文字连续粘贴。

同样，ROM/RAM 的零基础定义、QEMU 128 KiB ROM、实体 SPI Flash 和模块 UEFI
是同一个词在四个边界上的用法。书中先给一般概念，再在每条具体链上收窄：

\[
\begin{aligned}
\text{非易失启动存储的一般概念}
&\rightarrow \text{QEMU 自研 ROM}\\
&\rightarrow \text{载板可选 SPI Flash}\\
&\rightarrow \text{实体模块 UEFI 平台契约}.
\end{aligned}
\]

这种融合方式既保留全部知识责任，又避免相互矛盾的重复段落。若从原分册反查，
上面的表给出主落点；若从书顺序学习，正文已经把前置概念放在第一次需要它之前。

\section*{六、读者怎样验证“已经融合”}

\begin{enumerate}[leftmargin=2.2em]
  \item 在本索引找原 learning 文件，定位对应章节；
  \item 在章节先找术语定义，再找状态表、图、失败框和验收证据；
  \item 对图中任一箭头回答源、目的、载荷、初始化者、完成条件和失败所有权；
  \item 对实体网络回到十页原理图和 KiCad，不把学习图当制造数据；
  \item 对实现数字回到当前源码/CTest，不把历史阶段值当当前版本；
  \item 在书末术语表回查跨层同名概念，再用验收清单确认能够独立复述完整链。
\end{enumerate}

书稿输入检查会递归验证全部章节和图片存在，图形几何检查会验证七张整机图的
卡片净距以及三张实体电路图的显式引脚/导线连通；最终 XeLaTeX 构建和抽页渲染
再检查横置页、图注、字体和边界。自动化证明的是资源与几何底线，概念正确性
仍以源码、处理器/设备规范和固定上游硬件资料交叉审查。
